\section{A Priori Evaluation}
\label{sec:apriori}

\subsection{Validation and test set accuracy}

Table~\ref{tab:apriori_rmse} reports the root-mean-square error (RMSE) on the validation and test sets for all models.
The TBNN and TBRF models predict the full $6$-component anisotropy tensor $b_{ij}$, while the MLP models predict a scalar proxy (anisotropy magnitude); RMSE values are therefore not directly comparable across model types but are consistent within each group.

\begin{table}[htbp]
\centering
\caption{A priori accuracy: validation and test set RMSE. TBNN/PI-TBNN/TBRF RMSE is computed over all 6 anisotropy components; MLP RMSE is for the scalar target. Best tensor-predicting model on each metric in bold. The test set comprises periodic hills at $\alpha=1.2$ (interpolation within the training geometry class) and curved backward-facing step at $\Rey=13{,}700$ (unseen geometry).}
\label{tab:apriori_rmse}
\begin{tabular}{lrrrrrr}
\toprule
Model & Val RMSE & Test RMSE & PH $\alpha{=}1.2$ & CBFS & Degradation & Epochs \\
\midrule
TBRF (200 trees) & \textbf{0.0637} & \textbf{0.1253} & \textbf{0.0585} & \textbf{0.1435} & $1.97\times$ & n/a \\
PI-TBNN          & 0.1203          & 0.1418          & 0.0999          & 0.1553          & $1.18\times$ & 729 \\
MLP-Large        & 0.1045          & 0.1772          & 0.0797          & 0.2034          & $1.70\times$ & 344 \\
MLP              & 0.1096          & 0.1843          & 0.0886          & 0.2106          & $1.68\times$ & 1{,}000 \\
TBNN             & 0.0845          & 0.3889          & 0.0755          & 0.4573          & $4.60\times$ & 694 \\
\bottomrule
\end{tabular}
\end{table}

Among tensor-predicting models, the TBRF (200 trees) achieves the lowest validation RMSE (0.064), followed by the TBNN (0.085, 33\% higher).
The 10-tree compact TBRF captures 96\% of the full forest's accuracy improvement over TBNN, suggesting diminishing returns beyond $\sim$10 trees.

The most striking result in Table~\ref{tab:apriori_rmse} is the \emph{reversal} of model rankings between validation and test sets.
The TBNN achieves the best validation RMSE among neural networks (0.085) but the worst test RMSE (0.389)---a $4.6\times$ degradation.
In contrast, the PI-TBNN degrades only $1.18\times$ (0.120 $\to$ 0.142), and the TBRF degrades $1.97\times$ (0.064 $\to$ 0.125).
This reversal is driven almost entirely by the CBFS case (unseen geometry): the TBNN's CBFS RMSE is 0.457, compared to 0.076 on PH $\alpha=1.2$ (an interpolation case within the periodic hills geometry class).

\subsubsection{Per-case breakdown}

Table~\ref{tab:apriori_percase} decomposes the test set RMSE by case.
All models perform well on PH $\alpha=1.2$, which interpolates between training cases at $\alpha=0.5$, $0.8$, $1.0$, and $1.5$.
The CBFS geometry, however, introduces a recirculation region with qualitatively different separation and reattachment physics, exposing each model's ability to extrapolate.

\begin{table}[htbp]
\centering
\caption{Per-case test set RMSE for tensor-predicting models. PH $\alpha=1.2$ is an interpolation case; CBFS $\Rey=13{,}700$ is an unseen geometry. The CBFS/PH ratio quantifies how much worse each model performs on the new geometry.}
\label{tab:apriori_percase}
\begin{tabular}{lrrr}
\toprule
Model & PH $\alpha{=}1.2$ & CBFS $\Rey{=}13{,}700$ & CBFS/PH ratio \\
\midrule
TBRF     & 0.0585 & 0.1435 & $2.45\times$ \\
PI-TBNN  & 0.0999 & 0.1553 & $1.55\times$ \\
TBNN     & 0.0755 & 0.4573 & $6.05\times$ \\
\bottomrule
\end{tabular}
\end{table}

The TBNN's CBFS/PH ratio of $6.05\times$ indicates catastrophic overfitting to the training geometries.
The PI-TBNN, despite sharing the same architecture and parameter count, achieves a CBFS/PH ratio of only $1.55\times$---the most geometry-robust neural network model.
The TBRF's $2.45\times$ ratio reflects the ensemble averaging inherent in random forests, which provides natural regularization against overfitting.

\subsection{Generalization and overfitting}
\label{sec:generalization}

The validation-to-test degradation patterns in Table~\ref{tab:apriori_rmse} reveal a fundamental accuracy--generalization tradeoff.
Figure~\ref{fig:val_test_gap} summarizes this tradeoff: the TBNN sits in the high-accuracy, poor-generalization corner, while the PI-TBNN sacrifices validation accuracy for dramatically better test-set robustness.

The TBNN's overfitting is component-specific.
On the validation set, the TBNN's component-wise RMSE values are uniformly low (e.g., $b_{12}$: 0.053, $b_{11}$: 0.147).
On the CBFS test case, $b_{11}$ RMSE explodes to 0.806 ($5.5\times$ degradation) and $b_{22}$ to 0.580 ($6.8\times$), while $b_{12}$ degrades more modestly to 0.279 ($5.3\times$).
This suggests the TBNN has memorized geometry-specific normal stress magnitudes without learning the underlying physics governing separation and reattachment.

The PI-TBNN's regularization mechanism is noteworthy.
On the validation set, the realizability penalty appears to hurt accuracy ($+42\%$ RMSE vs TBNN).
However, this penalty constrains the network weights to produce outputs near the realizability boundary, limiting the effective capacity of the network and preventing it from fitting geometry-specific artifacts in the training data.
The result is a network that generalizes far better: PI-TBNN's CBFS RMSE (0.155) is $2.9\times$ better than the TBNN's (0.457).
This is a classic bias--variance tradeoff: the PI-TBNN accepts higher bias (worse validation fit) for dramatically lower variance (better generalization).

\subsection{Component-wise accuracy}

Table~\ref{tab:component_rmse} reports the component-wise RMSE on validation and test sets for all tensor-predicting models.

\begin{table}[htbp]
\centering
\caption{Component-wise RMSE on the test set for tensor-predicting models. Values in parentheses are the validation set RMSE for comparison. Components $b_{13}$ and $b_{23}$ are near-zero in the 2D-mean flows and are omitted.}
\label{tab:component_rmse}
\begin{tabular}{lrrr}
\toprule
Component & TBRF & TBNN & PI-TBNN \\
\midrule
$b_{11}$ & 0.167 (0.108) & 0.684 (0.147) & 0.203 (0.220) \\
$b_{12}$ & 0.071 (0.043) & 0.238 (0.053) & 0.074 (0.053) \\
$b_{22}$ & 0.129 (0.066) & 0.492 (0.085) & 0.172 (0.136) \\
$b_{33}$ & 0.211 (0.078) & 0.373 (0.104) & 0.211 (0.128) \\
\bottomrule
\end{tabular}
\end{table}

Several patterns emerge.
First, the shear stress component $b_{12}$ is predicted most accurately by all models, both on validation and test sets, consistent with its direct connection to the mean velocity gradient via the Boussinesq hypothesis.
Second, the TBNN's test-set degradation is concentrated in the normal stress components ($b_{11}$, $b_{22}$, $b_{33}$), where it shows $4$--$7\times$ degradation from validation to test.
Third, the PI-TBNN's $b_{12}$ test RMSE (0.074) is comparable to the TBRF (0.071) and far better than the TBNN (0.238), demonstrating that the realizability penalty preserves shear stress accuracy while preventing normal stress overfitting.

\subsection{Scatter plots}

\todo{Add predicted vs.\ true $b_{ij}$ scatter plots on the test set for key components ($b_{12}$, $b_{11}$) across all tensor-predicting models. Include $R^2$ values.}

\subsection{Lumley triangle}

\todo{Map predicted anisotropy eigenvalues onto the Lumley triangle \citep{Lumley1978} for each model. This provides a visual check of realizability: predictions outside the triangle violate physical constraints on the Reynolds stress tensor.}

\subsection{PI-TBNN beta sweep}
\label{sec:pi_sweep}

The realizability penalty weight $\beta$ in Eq.~\eqref{eq:pi_loss} was swept to determine whether physics-informed constraints improve accuracy.
Table~\ref{tab:pi_sweep} summarizes the results.

\begin{table}[htbp]
\centering
\caption{PI-TBNN realizability penalty sweep. Validation RMSE is pure MSE (no penalty terms).}
\label{tab:pi_sweep}
\begin{tabular}{rrrll}
\toprule
$\beta$ & Val RMSE & Epochs & $\Delta$ vs TBNN & Observation \\
\midrule
0 (TBNN)   & 0.0845 & 694 & ---      & Baseline \\
0.001      & 0.0852 & 729 & $+0.8\%$  & Negligible difference \\
0.01       & 0.0909 & 676 & $+7.6\%$  & Penalty degrades accuracy \\
\bottomrule
\end{tabular}
\end{table}

The unconstrained TBNN already produces near-realizable outputs.
Post-hoc analysis of TBNN predictions on the validation set shows that only 1.29\% of points (309 out of 23{,}967) violate the $b_{ii} \geq -1/3$ bound, and only 0.004\% violate $b_{ii} \leq 2/3$.
The worst violation is $b_{ii} = -0.45$ versus the limit of $-0.33$, and the total realizability penalty magnitude is 0.25\% of the MSE loss.

Based on validation accuracy alone, the penalty appears harmful---consistent with the findings of \citet{Wu2018} that loss-based penalties are redundant given architectural constraints.
However, the test set results (Table~\ref{tab:apriori_rmse}) dramatically reverse this conclusion: the PI-TBNN's test RMSE (0.142) is $2.7\times$ better than the TBNN's (0.389), despite its higher validation RMSE.
The realizability penalty acts as a regularizer that constrains the network's effective capacity, preventing overfitting to geometry-specific features in the training data (see Section~\ref{sec:generalization} for detailed analysis).

\subsection{TBRF tree count sweep}

Table~\ref{tab:tbrf_sweep} shows how TBRF accuracy varies with the number of trees per coefficient.

\begin{table}[htbp]
\centering
\caption{TBRF compact variants: accuracy vs.\ model size.}
\label{tab:tbrf_sweep}
\begin{tabular}{rrrrr}
\toprule
Trees & Total nodes & Binary size & Val RMSE & vs.\ TBNN \\
\midrule
1   & 282{,}902    & \SI{5.7}{MB}   & 0.0778 & $7.9\%$ better \\
5   & 1{,}432{,}612 & \SI{28.7}{MB}  & 0.0678 & $19.8\%$ better \\
10  & 2{,}797{,}130 & \SI{55.9}{MB}  & 0.0650 & $23.1\%$ better \\
200 & 55{,}051{,}026 & \SI{1{,}101}{MB} & 0.0637 & $24.6\%$ better \\
\bottomrule
\end{tabular}
\end{table}

The accuracy improvement saturates rapidly: 10 trees capture 96\% of the full 200-tree ensemble's improvement over TBNN, while requiring only 5\% of the model size.
However, even the 1-tree variant (\SI{5.7}{MB}) involves branching tree traversals that are poorly suited for GPU execution, compared to the TBNN's \SI{196}{KB} of dense matrix multiplies.
This tradeoff between offline accuracy and online deployability is central to the cost analysis in Section~\ref{sec:cost}.
